%% 2i2d-pont-thermique — liaison plancher/façade : isolation intérieure (ITI) contre extérieure (ITE).
%% Même géométrie dans les deux vignettes : mur vertical + plancher horizontal en béton.
%% ITI : l'isolant est interrompu par le plancher, les lignes de flux se resserrent -> point froid.
%% ITE : l'isolant est continu, les lignes de flux restent parallèles et espacées.
\documentclass[tikz,border=4pt]{standalone}
\usepackage{liaisons}
\usetikzlibrary{decorations.pathmorphing}
\begin{document}
\begin{tikzpicture}[x=1cm,y=1cm,
  txt/.style={font=\scriptsize},
  amb/.style={font=\tiny\bfseries},
  flux/.style={-{Stealth[length=4pt]}, line width=0.9pt, solideD},
  beton/.style={fill=black!18, draw=black!70, line width=0.6pt}]

\def\hv{2.4}     % hauteur d'une vignette
\def\dx{3.95}    % décalage de la vignette de droite

%% ---------- béton : mur vertical (0,9 -> 1,35) et plancher (1,35 -> 2,9) ---
\newcommand\betons{%
  \draw[beton] (0.90,0) rectangle (1.35,\hv);
  \draw[beton] (1.35,1.05) rectangle (2.90,1.35);}
%% ---------- zigzag « isolant » dans un rectangle ---------------------------
\newcommand\isolant[4]{%
  \fill[solideE!30] (#1,#2) rectangle (#3,#4);
  \begin{scope}
    \path[clip] (#1,#2) rectangle (#3,#4);
    \foreach \k in {0,1,...,20} {
      \pgfmathsetmacro\yy{#2+0.13*\k}
      \ifdim\yy cm<#4 cm
        \draw[solideE!80!black, line width=0.45pt, decorate,
              decoration={zigzag, amplitude=1.1pt, segment length=2.8pt}]
          (#1,\yy) -- (#3,\yy);
      \fi }
  \end{scope}
  \draw[line width=0.5pt, black!70] (#1,#2) rectangle (#3,#4);}

%% ======================= vignette gauche : ITI ==============================
\begin{scope}
  \fill[solideB!12] (0,0) rectangle (0.90,\hv);
  \fill[solideA!12] (1.60,0) rectangle (3.30,\hv);
  \fill[solideA!12] (1.35,1.05) rectangle (1.60,1.35);
  \isolant{1.35}{1.35}{1.60}{\hv}          % isolant au-dessus du plancher
  \isolant{1.35}{0}{1.60}{1.05}            % isolant au-dessous : interrompu
  \betons
  %% faisceau resserré : toutes les lignes contournent l'isolant par le plancher
  \foreach \y in {0.35, 0.80, 1.20, 1.70, 2.10} {
    \pgfmathsetmacro\ye{1.20+0.30*(\y-1.20)}
    \draw[flux] (2.95,\y) .. controls (2.10,\y) and (1.95,1.20) .. (1.35,1.20)
                          .. controls (0.95,1.20) and (0.70,\ye) .. (0.10,\ye); }
  \fill[solideA] (1.35,1.20) circle (0.085);
  \draw[line width=0.4pt, black!70] (1.42,1.28) -- (1.62,1.72);
  \node[txt, text=solideA, fill=white, inner sep=1pt, anchor=south west] at (1.55,1.70) {point froid};
  \node[amb, text=solideB, fill=solideB!12, inner sep=1pt] at (0.45,0.18) {EXT.};
  \node[amb, text=solideA, fill=solideA!12, inner sep=1pt] at (2.85,2.22) {INT.};
  \draw[line width=0.6pt, black!55] (0,0) rectangle (3.30,\hv);
  \node[font=\scriptsize\bfseries, anchor=south, align=center] at (1.65,\hv+0.10)
    {Isolation par\\ l'intérieur (ITI)};
  \node[txt, anchor=north, text=solideD] at (1.65,-0.10) {$\psi = 0{,}60$ W/(m$\cdot$K)};
\end{scope}

%% ======================= vignette droite : ITE ==============================
\begin{scope}[shift={(\dx,0)}]
  \fill[solideB!12] (0,0) rectangle (0.65,\hv);
  \fill[solideA!12] (1.35,0) rectangle (3.30,\hv);
  \isolant{0.65}{0}{0.90}{\hv}             % isolant continu, côté extérieur
  \betons
  %% lignes de flux parallèles et espacées : pas de resserrement
  \foreach \y in {0.35, 0.80, 1.20, 1.70, 2.10} {
    \draw[flux] (2.95,\y) -- (0.10,\y); }
  \node[amb, text=solideB, fill=solideB!12, inner sep=1pt] at (0.33,0.18) {EXT.};
  \node[amb, text=solideA, fill=solideA!12, inner sep=1pt] at (2.85,2.22) {INT.};
  \draw[line width=0.6pt, black!55] (0,0) rectangle (3.30,\hv);
  \node[font=\scriptsize\bfseries, anchor=south, align=center] at (1.65,\hv+0.10)
    {Isolation par\\ l'extérieur (ITE)};
  \node[txt, anchor=north, text=solideE!60!black] at (1.65,-0.10) {$\psi$ fortement réduit};
\end{scope}

%% ======================= relation commune ===================================
\node[draw=black!70, line width=0.7pt, rounded corners=2pt, fill=black!3, align=center,
      inner sep=4pt, anchor=north] at (3.62,-0.62)
  {{\small $\Phi = \psi \times L \times \Delta T$}\\[1pt]
   {\scriptsize $\Phi$ en W ; $\psi$ en W/(m$\cdot$K) ; $L$ en m ; $\Delta T$ en K}};
\end{tikzpicture}
\end{document}
